Un cliente detecta un fallo cuando se vence el tiempo de espera, cuando llega un 5xx o cuando llega un SOAP Fault. Para recuperarse puede reintentar con esperas cada vez más largas, pero solo si el método es idempotente, porque repetir un \texttt{POST} puede crear dos veces lo mismo \cite{rfc9110}. También puede tener un plan B, como hace SGIRI.

El caso real es la caída de AWS del 19 y 20 de octubre de 2025, en horario del Pacífico. A las 23:48 empezaron los errores en la API de DynamoDB; a las 00:38 los ingenieros ubicaron el problema en el DNS; a las 02:25 lo restauraron y a las 02:40 DynamoDB volvió a la normalidad. Fueron 2 horas y 52 minutos, pero el efecto en EC2, NLB y Lambda duró hasta las 14:20, en total 14 horas y 32 minutos \cite{aws2025}. La causa: dos procesos de la automatización del DNS chocaron, uno aplicó un plan viejo encima del nuevo y la limpieza lo borró, así que el nombre del servicio quedó sin direcciones IP. En el Ishikawa ordenamos las causas en automatización, procesos, dependencias y recuperación, con el mismo método de AC-01 \cite{soldani2023}. AWS apagó esa automatización en todo el mundo hasta corregirla.
